\documentclass{beamer}
\mode<presentation> {\usetheme{default}}
%\setbeamersize{text margin left=10mm,text margin right=10mm}


\setbeamertemplate{footline}[frame number]

\usepackage{amssymb}
\usepackage{amsthm}
\usepackage{amsmath}
\usepackage{booktabs}
\usepackage{needspace}
\usepackage{comment}
\usepackage{tikz}
\usepackage{graphicx}
\usepackage{bbm, dsfont}
\usepackage{array}

\usetikzlibrary{arrows,shapes,calc,positioning}

\theoremstyle{plain}
\newtheorem{thm}{Theorem}[section]
\newtheorem{cor}[thm]{Corollary}
\newtheorem{lem}[thm]{Lemma}
\newtheorem{prop}[thm]{Proposition}
%
\theoremstyle{definition}
\newtheorem{defn}{Definition}[section]
\newtheorem{remk}{Definition}[section]
\newtheorem{exam}{Example}[section]
\newtheorem{prob}{Exercise}[section]
\newtheorem{solu2}{Solution}[section]
\newenvironment{Shaded}{\begin{snugshade}}{\end{snugshade}}

\newcommand{\dif}{\operatorname{d}}
\newcommand{\var}{\operatorname{Var}}
\newcommand{\blue}[1]{\textcolor{blue}{#1}}
\newcommand{\purple}[1]{\textcolor{purple}{#1}}

% Chapter 10 has no numbered sections
\renewcommand{\thesection}{10.\arabic{section}}

\definecolor{accent}{RGB}{237,125,49}
\definecolor{dgreen}{RGB}{84,130,53}
\newcommand{\orange}[1]{\textcolor{accent}{#1}}
\newcommand{\green}[1]{\textcolor{dgreen}{#1}}

% Simple timeline: \timeline{length}{list of position/label pairs}
\newcommand{\timeline}[2]{%
  \begin{tikzpicture}
    \draw (0,0) -- (#1,0);
    \foreach \x/\lab in {#2} {
      \draw (\x,0.12) -- (\x,-0.12);
      \node[below] at (\x,-0.12) {\scriptsize $\lab$};
    }
  \end{tikzpicture}%
}

% Timeline 0, 1, 2 with row labels underneath:
% \swaptimeline{comma-separated row labels}
\newcommand{\swaptimeline}[1]{%
  \begin{tikzpicture}
    \draw (3.5,0) -- (9.5,0);
    \foreach \x/\lab in {4/0, 6.5/1, 9/2} {
      \draw (\x,0.12) -- (\x,-0.12);
      \node[above] at (\x,0.12) {\scriptsize $\lab$};
    }
    \foreach \lab [count=\k] in {#1} {
      \node[anchor=west] at (0,-0.6*\k) {\scriptsize \lab};
    }
  \end{tikzpicture}%
}

\title{RISK 201: Interest Rate Swaps}
\author{Himchan Jeong, FSA}
\institute{Department of Statistics \& Actuarial Science \\Simon Fraser University}
\date{}

\begin{document}
\maketitle

\begin{frame}{Fixed vs variable interest rate}

\begin{itemize}
\item A \purple{fixed interest rate}:
  \begin{itemize}
  \item is determined at the time the loan is written,
  \item does not change over the term of the loan.
  \end{itemize}
\medskip
\item A \purple{variable interest rate}:
  \begin{itemize}
  \item is adjusted periodically to reflect market conditions,
  \item is usually expressed as an index (e.g., the prime rate, or the London Inter-Bank Offered Rate) plus a spread.
  \end{itemize}
\end{itemize}

\end{frame}

\begin{frame}{Interest rate swaps}

An interest rate swap is a contractual agreement
\begin{itemize}
\item between two parties (the \orange{counterparties}) in which
\item each party makes periodic interest payments to the other
\item over a specified period (the \orange{swap term} or \orange{swap tenor})
\item based on a specific principal amount (the \orange{notional amount}).
\end{itemize}

\end{frame}

\begin{frame}{The transaction}

\begin{itemize}
\item One party (the \orange{payer}) pays the \orange{swap rate} (fixed).
\item The other party (the \orange{receiver}) pays a variable rate.
\end{itemize}

\vfill
\end{frame}

\begin{frame}{Timing of payments}

\begin{itemize}
\item The specific times at which payments are made (the \orange{settlement dates}) are usually equally spaced (daily, weekly, monthly, quarterly, annually).
\item The time between settlement dates is known as the \orange{settlement period}.
\item The first settlement period usually begins at time 0.
\item A swap that starts at a later date is a \orange{deferred swap}.
\end{itemize}

\end{frame}

\begin{frame}{Notional amounts, $Q_t$}

\begin{itemize}
\item Most interest rate swaps have a \orange{level notional amount} over the swap term: $Q_t = Q$.
\medskip
\item A swap with a varying notional amount that follows the outstanding balance of an amortization loan is known as an \orange{amortizing swap}.
\medskip
\item A swap with an increasing notional amount is an \orange{accreting swap}.
\end{itemize}

\end{frame}

\begin{frame}{Uses of interest rate swaps}

\begin{itemize}
\item \green{Hedging} $\rightarrow$ reducing risk associated with an existing obligation
  \begin{itemize}
  \item If a company has an obligation to pay a variable rate (e.g., on a loan) but would prefer to pay a fixed rate (e.g., because their income is fixed)
  \item If a company has an obligation to pay a fixed rate but would prefer to pay a variable rate (e.g., because income is interest-sensitive)
  \end{itemize}
\item \textcolor{red}{Speculation} $\rightarrow$ making a bet in search of a reward
\end{itemize}

\end{frame}

\begin{frame}{Example 10.1}

ABC Corp has borrowed \$500,000 for the next 2 years at a variable interest rate. Under this loan, ABC will pay interest at the end of year 1 and the end of year 2. The amount of the interest is based on the one-year spot rate at the start of year 1 and year 2, respectively.

\vspace{2em}
\begin{center}
\timeline{6}{0/0, 3/1, 6/2}
\end{center}

\vfill
\end{frame}

\begin{frame}{Example 10.1, continued}

\begin{enumerate}[a)]
\item Suppose the current spot rates are an annual effective interest rate of 5\% for a one-year period and an annual effective interest rate of 6\% for a two-year period. Find the implied one-year forward rate.
\end{enumerate}

\vspace{2em}
\begin{center}
\timeline{6}{0/0, 3/1, 6/2}
\end{center}

\vfill
\end{frame}

\begin{frame}{Example 10.1, continued}

\begin{enumerate}[a)]
\setcounter{enumi}{1}
\item ABC wants to enter into an interest rate swap that will fix the interest rate for two years. Under the swap, ABC will pay a fixed rate $R$ during both years of the loan. Find the swap rate $R$.
\end{enumerate}

\vspace{2em}
\begin{center}
\timeline{6}{0/0, 3/1, 6/2}
\end{center}

\vfill
\end{frame}

\begin{frame}{General formula for swap rate}

\begin{itemize}
\item Consider a swap with annual settlement and notional amounts $Q_t$.
\item The swap rate is such that, at issue, the PV of the payments on the fixed leg = PV of payments on variable leg
\end{itemize}

\begin{center}
\begin{tikzpicture}
  \draw (3.5,0) -- (6.3,0);
  \node at (6.9,0) {$\cdots$};
  \draw (7.5,0) -- (10,0);
  \foreach \x/\lab in {4/0, 5/1, 6/2, 9.5/n} {
    \draw (\x,0.12) -- (\x,-0.12);
    \node[above] at (\x,0.12) {\scriptsize $\lab$};
  }
  \node[anchor=west] at (-0.5,-0.8) {\scriptsize Fixed rate payments:};
  \node[anchor=west] at (-0.5,-2.0) {\scriptsize Variable rate payments:};
\end{tikzpicture}
\end{center}

\vfill
\end{frame}

\begin{frame}{General formula for swap rate, cont'd}

\[
\sum_{t=1}^{n} [Q_t \cdot R]\,(1+s_t)^{-t} = \sum_{t=1}^{n} [Q_t \cdot {}_{t-1}i_1]\,(1+s_t)^{-t}
\]

Let $P_t = (1+s_t)^{-t}$ denote the price of a $t$-year zero-coupon bond.
\[
\sum_{t=1}^{n} [Q_t \cdot R]\,P_t = \sum_{t=1}^{n} [Q_t \cdot {}_{t-1}i_1]\,P_t
\]
\[
\Rightarrow R = \frac{\purple{\sum_{t=1}^{n} [Q_t \cdot {}_{t-1}i_1]\,P_t}}{\sum_{t=1}^{n} Q_t P_t}
\]

\end{frame}

\begin{frame}{Example 10.2}

Consider an accreting swap with a 3-year term and annual settlement periods, with $Q_1 = 1000$, $Q_2 = 2000$ and $Q_3 = 3000$. You are given the following spot rates: $s_1 = 1.35\%$, $s_2 = 2.40\%$, $s_3 = 3.50\%$.

\medskip
Find the swap rate.

\vspace{1em}
\begin{center}
\timeline{6}{0/0, 2/1, 4/2, 6/3}
\end{center}

\vfill
\end{frame}

\begin{frame}{Level notional amount}

\textbf{$Q_t$ is constant $\rightarrow$ cancel from numerator/denominator}

\medskip
Also, $\displaystyle {}_{t-1}i_1 = \frac{(1+s_t)^t}{(1+s_{t-1})^{t-1}} - 1 = \frac{P_{t-1}}{P_t} - 1.$

\[
R = \frac{\sum_{t=1}^{n} \left[\frac{P_{t-1}}{P_t} - 1\right] P_t}{\sum_{t=1}^{n} P_t}
  = \frac{\sum_{t=1}^{n} [P_{t-1} - P_t]}{\sum_{t=1}^{n} P_t}
  = \frac{P_0 - P_n}{\sum_{t=1}^{n} P_t}
\]
\[
R = \frac{1 - P_n}{\sum_{t=1}^{n} P_t}
\]

\end{frame}

\begin{frame}{Example 10.1, revisited}

\begin{itemize}
\item $Q_t = \$500{,}000$ for $t = 1, 2$.
\item $s_1 = 5\%$, $s_2 = 6\%$
\end{itemize}

\vfill
\end{frame}

\begin{frame}{Net swap payments}

Under an interest rate swap, the payments are \orange{netted}: only the difference between the payments (the \orange{net swap payment}) is made by one party to the other.

\medskip
\underline{Example 10.1, continued}: Suppose ABC enters into the swap with DEF Ltd. Find the net swap payments.

\begin{center}
\swaptimeline{Fixed rate payment:, Variable rate payment:, Net swap payment:}
\end{center}

\vfill
\end{frame}

\begin{frame}{Market value of a swap}

\begin{itemize}
\item The \orange{market value} of an interest rate swap is the present value of the expected future cash flows under the agreement (ignore transaction costs).
\item The PV is determined using the spot rates applicable at the time the market value is calculated.
\item The MV of the swap at $t = 0$ is zero, since PV of the fixed rate pmts = PV of the variable rate pmts.
\item The MV \orange{after} the swap commences can be positive or negative.
\end{itemize}

\end{frame}

\begin{frame}{Example 10.1, continued}

$Q = \$500{,}000$, $s_1 = 5\%$, $s_2 = 6\%$, $R = 5.97074\%$

Assuming one year has passed, determine the market value of the swap if:
\begin{enumerate}[a)]
\item the one-year spot rate at time 1 is 7.0095238\%.
\end{enumerate}

\begin{center}
\swaptimeline{Fixed rate payment:, Variable rate payment:}
\end{center}

\vfill
\end{frame}

\begin{frame}{Example 10.1, continued}

\begin{enumerate}[a)]
\setcounter{enumi}{1}
\item the one-year spot rate at time 1 is 5\%.
\end{enumerate}

\begin{center}
\swaptimeline{Fixed rate payment:, Variable rate payment:}
\end{center}

\vfill
\end{frame}

\end{document}
